\documentclass[french]{article}
\usepackage{babel}
\def\frenchcontentsname{Sommaire}
\usepackage{lipsum}
\makeatletter
\def\maketitle{%
    \begin{center}\leavevmode
        \normalfont
        \rule[0pt]{\textwidth}{1pt}\par
        {\LARGE \@title\par}%
        {\small \@author\par}%
        {\Large \@date\par}%
        \rule[0pt]{\textwidth}{1pt}\par
    \end{center}%
}
\makeatother
\title{Code source}
\author{Maxime LE BLAN}
\date{}

%Choix du thème de la police (sans serif, petites capitales dans les titres...) 
\usepackage{amssymb} 
% Des symboles mathématiques
\usepackage{graphicx}
% Pour insérer des images
\usepackage{multicol}
% Pour gérer facilement un environnement à n colonnes
\usepackage{booktabs}
% Pour des beaux tableaux
\usepackage{esvect} 
% Pour de belles flèches sur les vecteurs
\usepackage{esdiff}
% Pour de beaux symboles de dérivation
\usepackage{siunitx}
% Pour l'utilisation des unités du système international
\sisetup{output-decimal-marker={,},inter-unit-product=\ensuremath{{}\cdot{}}}
% Quelques réglages pour l'affichage des unités

% Pour régler les marges de l'article
\usepackage{geometry}
\geometry{hmargin=2cm,vmargin=1.7cm}
% Pour les couleurs
\usepackage{color}
% Pour les indentations
\usepackage{tcolorbox}
\tcbuselibrary{skins, breakable}
\newtcolorbox{indentation}[1][]{
	blanker, left=1cm, borderline west={0.5mm}{0.6cm}{black!50}, breakable,
	before skip=2ex plus 0.1ex, after skip=2ex plus 0.1ex, #1
}
% Commande pour utiliser une indentation :
% \begin{indentation}
%     <texte>
% \end{indentation}

% Pour écrire du code en Python :
\usepackage[utf8]{inputenc}
\usepackage[T1]{fontenc}
\usepackage{xcolor}
\usepackage{listings}

\lstset{%
  language=Python,
  basicstyle   = \ttfamily,
  keywordstyle =    \color{red!70!gray},
  keywordstyle = [2]\color{blue},
  commentstyle =    \color{gray}\itshape,
  stringstyle  = \color{green!50!black},
  numberstyle = \color{gray},
  numbers = left,
  frame = single,
  framesep = 2pt,
  breaklines = true,
  aboveskip = 1ex,
  xleftmargin = 0mm,
  xrightmargin = 0mm,
  framexleftmargin = 8mm,
  literate        =
                        {é}{{\'e}}{1}%
                        {è}{{\`e}}{1}%
                        {à}{{\`a}}{1}%
                        {â}{{\^a}}{1}%%%
                        {ç}{{\c{c}}}{1}%
                        {œ}{{\oe}}{1}%
                        {ù}{{\`u}}{1}%
                        {É}{{\'E}}{1}%
                        {È}{{\`E}}{1}%
                        {À}{{\`A}}{1}%
                        {Ç}{{\c{C}}}{1}%
                        {Œ}{{\OE}}{1}%
                        {Ê}{{\^E}}{1}%
                        {ê}{{\^e}}{1}%
                        {î}{{\^i}}{1}%
                        {ï}{{\"i}}{1}%%%
                        {ô}{{\^o}}{1}%
                        {û}{{\^u}}{1}%
}

\begin{document}

\maketitle
\tableofcontents
\newpage

\section{Modules}
\lstinputlisting{modules.py}
Diapositive 36

\section{algos\_chiffrement.py}
\lstinputlisting{algos_1.py}
Diapositive 37 \\

\lstinputlisting{algos_2.py}
Diapositive 38 \\

\lstinputlisting{algos_3.py}
Diapositive 39 \\

\lstinputlisting{algos_4.py}
Diapositive 40 \\

\lstinputlisting{algos_5.py}
Diapositive 41 \\

\lstinputlisting{algos_6.py}
Diapositive 42 \\

\lstinputlisting{algos_7.py}
Diapositive 43 \\

\lstinputlisting{algos_8.py}
Diapositive 44 \\

\lstinputlisting{algos_9.py}
Diapositive 45 \\

\lstinputlisting{algos_10.py}
Diapositive 46 \\

\lstinputlisting{algos_11.py}
Diapositive 47 \\

\lstinputlisting{algos_12.py}
Diapositive 48

\section{serveur.py}
\lstinputlisting{serveur_1.py}
Diapositive 49 \\

\lstinputlisting{serveur_2.py}
Diapositive 50 \\

\lstinputlisting{serveur_3.py}
Diapositive 51 \\

\lstinputlisting{serveur_4.py}
Diapositive 52 \\

\lstinputlisting{serveur_5.py}
Diapositive 53 \\

\lstinputlisting{serveur_6.py}
Diapositive 54

\section{client.py}
\lstinputlisting{client_1.py}
Diapositive 55 \\

\lstinputlisting{client_2.py}
Diapositive 56 \\

\lstinputlisting{client_3.py}
Diapositive 57 \\

\lstinputlisting{client_4.py}
Diapositive 58 \\

\lstinputlisting{client_5.py}
Diapositive 59 \\

\lstinputlisting{client_6.py}
Diapositive 60 \\

\lstinputlisting{client_7.py}
Diapositive 61 \\

\lstinputlisting{client_8.py}
Diapositive 62 \\

\lstinputlisting{client_9.py}
Diapositive 63 \\

\lstinputlisting{client_10.py}
Diapositive 64 \\

\lstinputlisting{client_11.py}
Diapositive 65 \\

\lstinputlisting{client_12.py}
Diapositive 66 \\

\lstinputlisting{client_13.py}
Diapositive 67 \\

\lstinputlisting{client_14.py}
Diapositive 68

\section{fonctions\_utiles.py}
\lstinputlisting{fu_1.py}
Diapositive 69 \\

\lstinputlisting{fu_2.py}
Diapositive 70 \\

\lstinputlisting{fu_3.py}
Diapositive 71 \\

\lstinputlisting{fu_4.py}
Diapositive 72 \\

\lstinputlisting{fu_5.py}
Diapositive 73 \\

\lstinputlisting{fu_6.py}
Diapositive 74 \\

\lstinputlisting{fu_7.py}
Diapositive 75 \\

\lstinputlisting{fu_8.py}
Diapositive 76 \\

\lstinputlisting{fu_9.py}
Diapositive 77 \\

\lstinputlisting{fu_10.py}
Diapositive 78

\end{document}
